\begin{statement}{C.2.UN.31}{Proved}{Nonnegative-Kodaira-dimension components have tightly specified cyclic covers}
Let $f:U\to\Delta$ be an unrestricted proper degeneration
of connected complex two-tori from a smooth threefold.
Let $Z$ be the smooth minimal model of the normalization
of one central component, with $\kappa(Z)\ge0$.
There is an actual generically finite component cover
$E\to S^\nu$, obtained by the normalized root base change
and resolution, with cyclic stabilizer $G$, such that
$E/G$ is properly birational to $S^\nu$. Its smooth
minimal model $E_0$ satisfies one of these descriptions:

\begin{itemize}
\item $E_0$ is a complex two-torus;
\item $E_0$ has an isotrivial elliptic fibration over a
curve $B$ with $g(B)\le1$, identity integral $H_1$
monodromy, and only positive multiples of smooth
elliptic fibres. In this case
$R^1q_*\mathcal O_{E_0}=\mathcal O_B$,
$p_g(E_0)=g(B)$, and $\chi(\mathcal O_{E_0})=0$.
\end{itemize}
In both cases $\chi(\mathcal O_{E_0})=0$, and the
action of $G$ descends holomorphically to $E_0$.

For the original minimal component the consequences are
\[
 p_g(Z)\le1,
 \qquad \chi(\mathcal O_Z)\in\{0,1,2\}.
                                                               \tag{C.CS.1}
\]
Every component with $\kappa(Z)=1$ has an actual
isotrivial elliptic Iitaka fibration on $Z$, over a
curve of genus at most one. An Enriques component
has an isotrivial elliptic fibration over $\mathbb P^1$.
A K3 component is either the cyclic torus quotient
of C.2.UN.24, with the exact orders and configurations
listed there, or has an isotrivial elliptic fibration
over $\mathbb P^1$ through the non-torus-cover branch.
The singular and multiple fibres on these original
quotient surfaces are retained by the quotient maps;
the smooth-multiple assertion concerns the specified
cover $E_0$.
\end{statement}

\begin{proof}
Take the actual component, cyclic stabilizer, and
proper birational quotient comparison in C.2.UN.16
and (C.CQ.12). Pluricanonical pullback through a
resolution of the generically finite graph gives
\[
             \kappa(E)\ge\kappa(Z)\ge0.
                                                               \tag{C.CS.2}
\]
The compact-component theorem C.2.UN.15 applies
also to this finite base-change family, so
$\kappa(E)\le1$. For the saturated limiting field
space, the generic-rank-two branch is covered
by C.2.UN.22. In rank one the first integral
has rational or elliptic general fibre. Rational
fibres would force every positive pluricanonical
section to vanish, by the exact adjunction
argument of C.2.UN.20, contradicting (C.CS.2).
Thus C.2.UN.29 applies. These two statements
give precisely the asserted torus or smooth-multiple
elliptic alternatives and coherent Euler characteristic
zero.

In the elliptic alternative the actual root-disc
invariant calculation of (C.RG.2)--(C.RG.13) gives
$R^1q_*\mathcal O=\mathcal O_B$ and
$p_g(E)=g(B)$. That calculation requires only
the fixed elliptic modulus, identity integral
monodromy and the smooth-multiple local charts,
so it applies to the elliptic alternative of
C.2.UN.22 as well. The independent compact-core
bound C.2.UN.28 gives $p_g(E)\le1$, and hence
$g(B)\le1$ in both branches.

All proper birational maps between the displayed
smooth surface models factor through actual
point blowups and contractions. On a nonruled
surface the minimal model is unique. Equivalently,
the exceptional curves of a minimal contraction
form a negative-definite configuration; its
$(-1)$-curves are disjoint, so their finite
$G$-orbits can be contracted equivariantly,
one stage at a time. This gives the holomorphic
action on $E_0$ and the exact comparison
\[
 \begin{array}{ccc}
 E&\longrightarrow&E_0\\
 \downarrow&&\downarrow\\
 E/G&\longrightarrow&E_0/G
 \end{array}
 \qquad E_0/G\dashrightarrow Z
       \quad\hbox{properly birational}.
                                                               \tag{C.CS.3}
\]
It is this action and quotient which are used
below.

Suppose first that $\kappa(Z)=1$. Then
(C.CS.2) and C.2.UN.15 give $\kappa(E_0)=1$,
so the torus alternative is impossible. The
elliptic fibration is the connected-fibre
Iitaka fibration. To specify the invariance,
the group acts on every space of pluricanonical
sections. An equivariant graph resolution of
a sufficiently large pluricanonical map and
its Stein factor therefore give a $G$-equivariant
map to $B$. The same Stein base is the one
of the preceding elliptic fibration: pluricanonical
ratios restrict constantly to its elliptic
fibres, and Kodaira dimension one makes the
induced map of bases finite. Taking the
connected-fibre Stein factor removes that
finite map.

The actual quotient fibration is
\[
                   E'/G\longrightarrow B/G,
                                                               \tag{C.CS.4}
\]
after this equivariant graph resolution.
For a general base point, let
$K=\ker(G\to\operatorname{Aut}B)$.
Its fibre is exactly $F/K$, as in C.2.UN.20,
and is connected. Every $F$ is isomorphic
to the same elliptic curve $A$.
The subgroup $K$ acts by translations:
a finite automorphism of an elliptic curve
with nontrivial linear part has a fixed
point, and its quotient is rational by
Riemann--Hurwitz. A rational general fibre
in (C.CS.4) would annihilate every positive
pluricanonical section, contrary to
$\kappa(Z)=1$. The translation points
belong to the finite subgroup $A[|K|]$,
so they are locally constant under the
marked product identifications. Therefore
the quotient general fibre is the one
fixed elliptic curve $A/K$.
The quotient base has genus at most one
by Riemann--Hurwitz and $g(B)\le1$.

This fibration is the Iitaka fibration
of $Z$. Its map is holomorphic on the
minimal surface. One can also check the
descent on an actual common resolution
$Z\xleftarrow{\rho}X\xrightarrow q B/G$:
for a smooth general elliptic fibre $F$,
\[
 0=K_X\cdot F
   =\rho^*K_Z\cdot F+\sum_j a_jE_j\cdot F,
                   \qquad a_j>0.
                                                               \tag{C.CS.5}
\]
On a minimal surface of Kodaira dimension
one the canonical class is nef; this is
the elliptic canonical-bundle formula
in the surface classification used in
C.3.1. Every term in (C.CS.5) is
nonnegative. Hence every exceptional
curve is vertical. Connected exceptional
fibres and $\rho_*\mathcal O_X=\mathcal O_Z$
then descend the actual proper graph to
a holomorphic map on $Z$. This proves
the claimed isotrivial Iitaka fibration.

Now let $Z$ be Enriques. Algebraic
dimension is preserved by the actual
generically finite component map, by
the symmetric-function field argument
in (C.K3.2)--(C.K3.3). Thus $E_0$ is
projective. If it has Kodaira dimension
one, the preceding intrinsic Iitaka
quotient argument again gives an
isotrivial elliptic fibration on a
birational model of $Z$. If it has
Kodaira dimension zero, the surface
classification and $\chi(\mathcal O_{E_0})=0$
make it Abelian or bielliptic.
The Abelian case is excluded by the
exact cyclic quotient theorem C.2.UN.23.
For the bielliptic case, use the actual
orbit fibration of its connected
automorphism group. In its finite free
product presentation $(A_1\times A_2)/H$,
$H$ acts by translations on $A_1$ and
has nontrivial linear action on $A_2$.
The invariant tangent fields are exactly
the one-dimensional translation space
of $A_1$. Thus the image of the compact
elliptic translation group $A_1$ is
the whole connected automorphism group,
and its orbit map is the explicit
elliptic fibration to $A_2/H\simeq\mathbb P^1$.
Its fibres are smooth elliptic multiples;
its integral monodromy is the identity,
since all actions on the fibre direction
are translations. Every automorphism
normalizes the connected automorphism
group, so preserves this orbit map and
acts on its connected-fibre base. It
is therefore $G$-equivariant. The quotient
(C.CS.4) is available again. Its
general fibre must be elliptic, since
a rational general fibre would force
the nonzero even pluricanonical form
of an Enriques surface to vanish.

In either case, (C.CS.5) descends the
elliptic map to the minimal Enriques
surface: now $K_Z$ is torsion, so its
intersection term is zero. The Leray
injection of $H^1$ of the base into
$H^1(Z,\mathcal O_Z)=0$ makes the
base $\mathbb P^1$. This proves the
Enriques assertion through the actual
quotient fibration.

For a K3 target, the nonzero canonical
form gives $p_g(E_0)\ge1$. If $E_0$
is a torus, C.2.UN.23--24 give exactly
the cyclic symplectic torus quotient
already stated. Otherwise it has an
elliptic fibration of the specified
fixed-modulus type. If $\kappa(E_0)=1$,
use the equivariant Iitaka fibration.
If $a(E_0)=1$, use the intrinsic
field identity $\mathcal M(E_0)=q^*\mathcal M(B)$
and the equivariant Stein factor in
(C.K3.4). The remaining projective
Kodaira-dimension-zero case would be
bielliptic, but has $p_g=0$ and is
excluded. Algebraic dimension zero
has only the torus alternative by
C.2.UN.16. Thus these are all cases.
Taking the quotient gives an isotrivial
elliptic fibration, since its general
fibre cannot be rational on a K3
birational model. Formula (C.CS.5)
with $K_Z=0$ descends it to $Z$,
and $H^1(Z,\mathcal O_Z)=0$ again
makes its base $\mathbb P^1$.

Finally C.2.UN.28 gives $p_g(Z)\le1$.
The compact surface classification
gives $\chi(\mathcal O_Z)\ge0$ when
$\kappa(Z)\ge0$: in Kodaira dimension
zero its values are the explicit
torus, bielliptic, Kodaira, K3 and
Enriques values; in Kodaira dimension
one the relatively minimal elliptic
surface has $K_Z^2=0$ and nonnegative
Euler number by Kodaira's fibre list,
so surface Riemann--Roch gives the
inequality. On the other hand
\[
        \chi(\mathcal O_Z)=1-h^1(Z,\mathcal O_Z)+p_g(Z)
                              \le2.
                                                               \tag{C.CS.6}
\]
It is an integer, giving (C.CS.1).
The later exact quotient theorem
C.2.UN.32 further identifies every K3
component, in all algebraic dimensions,
with the cyclic torus quotient family.
The preceding non-torus-cover comparison
is retained as its explicit intermediate
fibration map.

All these comparisons retain the
finite quotient's singular fibres;
they do not identify them with the
smooth fibres of its root cover.
\end{proof}


% C.2.CR.1 from part_C2_contracted.tex

The statements C.2.CR.1--C.2.CR.6 below retain the
proved local and global reductions. The later
C.2.UN.7 proves that the contracted irrational
configuration has $e(S)\ge2$ and $\pi_1(S)=1$.
Together with the CDP transgression, or the zero
Euler conclusion C.2.MH.3, this excludes the
configuration entirely under the CDP hypotheses.
The unrestricted genus assertion is separately
proved in C.2.UN.10.

